\section{مدول‌های نوتری و آرتینی: ادامه}

% Source: sources/lectures/11 Notherian and Artinian modules (2).pdf | PDF page: 1 | item: finite generation/cogeneration examples
\begin{frame}[t]{تولید و هم‌تولید متناهی}
    \label{lec11:fg-fc-examples}
    \begin{example}
        $\mathbb Z$ به‌عنوان $\mathbb Z$-مدول با تولید متناهی است،
        ولی متناهیاً هم‌تولیدشده نیست.
        \[
            \bigcap_{\substack{p\ \text{اول}}}p\mathbb Z=0.
        \]
        \textbf{پرسشِ جزوه:} این اشتراک چگونه ادعای دوم را نشان می‌دهد؟
    \end{example}
    \begin{example}
        $\mathbb Z_{p^\infty}$ متناهیاً هم‌تولیدشده است،
        ولی با تولید متناهی نیست.
    \end{example}
\end{frame}

% Source: sources/lectures/11 Notherian and Artinian modules (2).pdf | PDF page: 1 | item: three equivalent Artinian conditions
\begin{frame}[t]{معیارهای آرتینی بودن}
    \label{lec11:artinian-criteria}
    \begin{theorem}
        برای هر $R$-مدول $M$، موارد زیر معادل‌اند:
        \begin{enumerate}
            \item $M$ آرتینی است.
            \item هر مدول خارج‌قسمتی $M$ متناهیاً هم‌تولیدشده است.
            \item هر مجموعهٔ ناتهی از زیرمدول‌های $M$ دارای عنصر مینیمال است.
        \end{enumerate}
    \end{theorem}
\end{frame}

% Source: sources/lectures/11 Notherian and Artinian modules (2).pdf | PDF page: 1 | item: Artinian implies finitely cogenerated quotients, initial setup
% AMBIGUITY[L11-Q01]: selection of finite intersections for arbitrary I is not explicit in the source; proof completion withheld.
\begin{frame}[t]{خارج‌قسمت‌های متناهیاً هم‌تولیدشده}
    \label{lec11:artinian-to-fc}
    \textbf{آغاز اثباتِ شرط ۱ به شرط ۲.}
    فرض کنید $M$ آرتینی و $N\leqslant M$ باشد. برای خانواده‌ای از زیرمدول‌های $M/N$،
    \[
        \bigcap_{i\in I}\frac{M_i}{N}=0
        \quad\Longrightarrow\quad
        \bigcap_{i\in I}M_i=N.
    \]
    باید زیرمجموعه‌ای متناهی $J\subseteq I$ پیدا کنیم که
    \[
        \bigcap_{i\in J}M_i=N.
    \]
    \begin{tblock}{یادداشت پیش‌نویس}
        انتخاب اشتراک‌های متناهی در ادامهٔ اثباتِ جزوه نیازمند روشن‌سازی است.
        این گام با شناسهٔ \lr{L11-Q01} صریحاً به بررسی بعدی موکول شده است.
    \end{tblock}
\end{frame}

% Source: sources/lectures/11 Notherian and Artinian modules (2).pdf | PDF page: 2 | item: finitely cogenerated quotients imply minimal condition, descending selection
\begin{frame}[t]{یافتن زیرمدول مینیمال}
    \label{lec11:fc-to-minimal-chain}
    \textbf{شرط ۲ به شرط ۳.}
    خانواده‌ای ناتهی از زیرمدول‌های $M$ را در نظر بگیرید و $M_1$ را از آن انتخاب کنید.
    اگر مینیمال باشد، کار تمام است؛ در غیر این صورت عضوی مانند
    $M_2\subsetneq M_1$ در همان خانواده وجود دارد.

    اگر این فرایند پایان نیابد، زنجیره‌ای به شکل
    \[
        M_1\supsetneq M_2\supsetneq M_3\supsetneq\cdots
    \]
    به‌دست می‌آید. قرار دهید
    \[
        N=\bigcap_{i\geqslant1}M_i\leqslant M.
    \]
\end{frame}

% Source: sources/lectures/11 Notherian and Artinian modules (2).pdf | PDF page: 2 | item: contradiction in quotient, minimal condition implies DCC
\begin{frame}[t]{پایان اثبات معیار آرتینی}
    \label{lec11:fc-to-minimal-end}
    در $M/N$ داریم
    \[
        \bigcap_{i\geqslant1}\frac{M_i}{N}=0.
    \]
    چون $M/N$ متناهیاً هم‌تولیدشده است و زنجیره تودرتو است، برای یک $k$ داریم
    \[
        \bigcap_{i=1}^{k}\frac{M_i}{N}=0
        \quad\Longrightarrow\quad M_k=N.
    \]
    در نتیجه $N=M_k=M_{k+1}=\cdots$؛ این با سره بودن شمول‌ها تناقض دارد.
    پس فرایند انتخاب پایان می‌یابد و یک عضو مینیمال به‌دست می‌آید.

    \medskip
    \textbf{شرط ۳ به شرط ۱:} در جزوه روشن دانسته شده است.
    \hfill\(\qed\)
\end{frame}

% Source: sources/lectures/11 Notherian and Artinian modules (2).pdf | PDF pages: 2-3 | item: Noetherian extension criterion and forward implication
\begin{frame}[t]{نوتری بودن زیرمدول و خارج‌قسمت}
    \label{lec11:noetherian-extension}
    \begin{theorem}
        اگر $N\leqslant M$، آنگاه $M$ نوتری است اگر و تنها اگر
        $N$ و $M/N$ نوتری باشند.
    \end{theorem}
    \textbf{اثباتِ جهت مستقیم.}
    هر زنجیرهٔ صعودی از زیرمدول‌های $N$، زنجیره‌ای از زیرمدول‌های $M$ نیز هست؛
    پس از نوتری بودنِ $M$ متوقف می‌شود.

    همچنین هر زنجیرهٔ صعودی در $M/N$ به شکل
    \[
        \frac{M_1}{N}\subseteq\frac{M_2}{N}\subseteq\frac{M_3}{N}\subseteq\cdots
    \]
    است. زنجیرهٔ متناظر در $M$ متوقف می‌شود و در نتیجه زنجیرهٔ خارج‌قسمت‌ها نیز متوقف می‌شود.
\end{frame}

% Source: sources/lectures/11 Notherian and Artinian modules (2).pdf | PDF page: 3 | item: converse Noetherian extension, quotient generators
\begin{frame}[t]{جهت عکس معیار نوتری: مولدهای خارج‌قسمت}
    \label{lec11:noetherian-extension-quotient}
    فرض کنید $N$ و $M/N$ نوتری باشند. ثابت می‌کنیم هر زیرمدول
    $K\leqslant M$ با تولید متناهی است.

    چون $(K+N)/N$ زیرمدولی از $M/N$ است، با تولید متناهی است.
    با قضیهٔ دوم یکریختی،
    \[
        \frac{K+N}{N}\cong\frac{K}{K\cap N}.
    \]
    بنابراین می‌توان نوشت
    \[
        \frac{K}{K\cap N}
        =\bigl\langle k_1+(K\cap N),\ldots,k_\ell+(K\cap N)\bigr\rangle,
    \]
    و در نتیجه
    \[
        K=\langle k_1,\ldots,k_\ell\rangle+(K\cap N).
    \]
\end{frame}

% Source: sources/lectures/11 Notherian and Artinian modules (2).pdf | PDF page: 3 | item: converse Noetherian extension, generators of intersection
\begin{frame}[t]{جهت عکس معیار نوتری: مولدهای زیرمدول}
    \label{lec11:noetherian-extension-generators}
    از $K\cap N\leqslant N$ و نوتری بودنِ $N$ نتیجه می‌شود
    \[
        K\cap N=\langle n_1,\ldots,n_m\rangle.
    \]
    پس
    \[
        K=\langle n_1,\ldots,n_m,k_1,\ldots,k_\ell\rangle.
    \]
    بنابراین هر زیرمدول $K$ با تولید متناهی است و $M$ نوتری است.
    \hfill\(\qed\)
\end{frame}

% Source: sources/lectures/11 Notherian and Artinian modules (2).pdf | PDF page: 3 | item: Artinian submodule/quotient direction
\begin{frame}[t]{آرتینی بودن زیرمدول و خارج‌قسمت}
    \label{lec11:artinian-inheritance}
    \begin{theorem}
        اگر $M$ آرتینی و $N\leqslant M$ باشد، $N$ و $M/N$ آرتینی‌اند.
    \end{theorem}
    \textbf{اثبات.} همان استدلال جهت مستقیمِ قضیهٔ نوتری،
    این بار برای زنجیره‌های نزولی، برقرار است.
    \hfill\(\qed\)

    \medskip
    عکس این قضیه نیز برقرار است. برای اثبات آن، ابتدا لم صفحهٔ بعد را در نظر می‌گیریم.
\end{frame}

% Source: sources/lectures/11 Notherian and Artinian modules (2).pdf | PDF page: 4 | item: modular lemma, statement
\begin{frame}[t]{لمی دربارهٔ جمع و اشتراک زیرمدول‌ها}
    \label{lec11:modular-lemma}
    \begin{lemma}
        فرض کنید $N_1,N_2,N_3$ زیرمدول‌های $M$ باشند و
        \[
            N_2\subseteq N_1,\qquad
            N_1+N_3=N_2+N_3,\qquad
            N_1\cap N_3=N_2\cap N_3.
        \]
        آنگاه $N_1=N_2$.
    \end{lemma}
\end{frame}

% Source: sources/lectures/11 Notherian and Artinian modules (2).pdf | PDF page: 4 | item: modular lemma, proof
\begin{frame}[t]{اثبات لم جمع و اشتراک}
    \label{lec11:modular-lemma-proof}
    با استفاده از فرض‌ها، مطابق محاسبهٔ جزوه داریم
    \[
        \begin{aligned}
            N_1
            &=N_1\cap(N_1+N_3)\\
            &=N_1\cap(N_2+N_3)\\
            &=(N_1\cap N_2)+(N_1\cap N_3)\\
            &=N_2+(N_2\cap N_3)\\
            &=N_2.
        \end{aligned}
    \]
    \hfill\(\qed\)
\end{frame}

% Source: sources/lectures/11 Notherian and Artinian modules (2).pdf | PDF pages: 4-5 | item: converse Artinian extension, induced chains
% AMBIGUITY[L11-Q02]: weak induced inclusions are a provisional correction of the strict source symbols.
% AMBIGUITY[L11-Q03]: Artinian replaces the two source occurrences of Noetherian provisionally; no user resolution is claimed.
\begin{frame}[t]{آرتینی بودن مدول از زیرمدول و خارج‌قسمت}
    \label{lec11:artinian-extension-pending}
    \begin{theorem}
        اگر $N\leqslant M$ و $N$ و $M/N$ آرتینی باشند، آنگاه $M$ آرتینی است.
    \end{theorem}
    \textbf{صورت موقت اثبات: \lr{L11-Q02} و \lr{L11-Q03}.}
    زنجیره‌ای نزولی از زیرمدول‌های $M$ را در نظر بگیرید. زنجیره‌های القایی عبارت‌اند از
    \[
        \begin{gathered}
            M_1+N\supseteq M_2+N\supseteq\cdots,\\
            \frac{M_1+N}{N}\supseteq\frac{M_2+N}{N}\supseteq\cdots,\\
            M_1\cap N\supseteq M_2\cap N\supseteq\cdots.
        \end{gathered}
    \]
    \textbf{اصلاح موقت:} شمول‌ها غیرسره نوشته شده‌اند؛ در جزوه سره‌اند.
    همچنین در ادامه از فرض آرتینی بودن استفاده می‌شود؛ جزوه در دو محل «نوتری» نوشته است.
\end{frame}

% Source: sources/lectures/11 Notherian and Artinian modules (2).pdf | PDF pages: 4-5 | item: stabilization of sum/quotient and intersection chains
% AMBIGUITY[L11-Q02], AMBIGUITY[L11-Q03]: source proof transcribed with explicitly provisional weak inclusions and Artinian hypotheses.
\begin{frame}[t]{توقف زنجیره‌های القایی: صورت موقت}
    \label{lec11:artinian-extension-stabilization}
    \textbf{ادامهٔ صورت موقت: \lr{L11-Q02} و \lr{L11-Q03}.}
    از آرتینی بودنِ $M/N$، شاخصی مانند $k$ وجود دارد که برای هر $i\geqslant0$،
    \[
        \frac{M_k+N}{N}=\frac{M_{k+i}+N}{N}
        \quad\Longrightarrow\quad M_k+N=M_{k+i}+N.
    \]
    همچنین از آرتینی بودنِ $N$، شاخصی مانند $t$ وجود دارد که برای هر $j\geqslant0$،
    \[
        M_t\cap N=M_{t+j}\cap N.
    \]
    بنابراین $k$ و $t$ به‌ترتیب شاخص‌های توقف زنجیرهٔ خارج‌قسمت‌ها و زنجیرهٔ اشتراک‌ها هستند.
\end{frame}

% Source: sources/lectures/11 Notherian and Artinian modules (2).pdf | PDF page: 5 | item: final modular-lemma step after stabilization
% AMBIGUITY[L11-Q02], AMBIGUITY[L11-Q03]: continuation of the explicitly provisional corrected proof.
\begin{frame}[t]{گام پایانی پس از توقف دو زنجیره}
    \label{lec11:artinian-extension-end}
    \textbf{ادامهٔ صورت موقت: \lr{L11-Q02} و \lr{L11-Q03}.}
    برای $s\geqslant\max\{k,t\}$ و هر $i\geqslant0$ داریم
    \[
        \begin{aligned}
            M_s+N&=M_{s+i}+N,\\
            M_s\cap N&=M_{s+i}\cap N,\\
            M_{s+i}&\subseteq M_s.
        \end{aligned}
    \]
    از لم قبل، $M_s=M_{s+i}$ و زنجیره متوقف می‌شود؛ در نتیجه $M$ آرتینی است.
    \hfill\(\qed\)

    \medskip
    اصلاح‌های این اثبات در پروندهٔ پرسش‌ها باز و ثبت‌شده باقی مانده‌اند.
\end{frame}

% Source: sources/lectures/11 Notherian and Artinian modules (2).pdf | PDF page: 5 | item: direct sum criterion, statement
\begin{frame}[t]{شرایط زنجیری برای جمع مستقیم}
    \label{lec11:direct-sum-criterion}
    \begin{theorem}
        اگر $\{M_i\}_{i\in I}$ خانواده‌ای از $R$-مدول‌های ناصفر باشد،
        آنگاه جمع مستقیم
        \[
            \bigoplus_{i\in I}M_i
        \]
        نوتری \textup{(آرتینی)} است اگر و تنها اگر $I$ متناهی باشد
        و هر $M_i$ نوتری \textup{(آرتینی)} باشد.
    \end{theorem}
\end{frame}

% Source: sources/lectures/11 Notherian and Artinian modules (2).pdf | PDF pages: 5-6 | item: infinite direct-sum chains
\begin{frame}[t]{چرا مجموعهٔ شاخص‌ها باید متناهی باشد؟}
    \label{lec11:direct-sum-infinite}
    اگر $I$ نامتناهی باشد، با حذف پی‌درپیِ شاخص‌های متمایز، زنجیرهٔ
    \[
        \bigoplus_{i\in I}M_i
        \supsetneq\bigoplus_{i\in I\setminus\{i_1\}}M_i
        \supsetneq\bigoplus_{i\in I\setminus\{i_1,i_2\}}M_i
        \supsetneq\cdots
    \]
    متوقف نمی‌شود؛ پس جمع مستقیم آرتینی نیست.

    همچنین زنجیرهٔ
    \[
        \{0\}\subsetneq M_{i_1}
        \subsetneq M_{i_1}\oplus M_{i_2}
        \subsetneq M_{i_1}\oplus M_{i_2}\oplus M_{i_3}
        \subsetneq\cdots
    \]
    متوقف نمی‌شود؛ پس جمع مستقیم نوتری نیست.
\end{frame}

% Source: sources/lectures/11 Notherian and Artinian modules (2).pdf | PDF page: 6 | item: finite direct-sum converse
\begin{frame}[t]{جمع مستقیم متناهی}
    \label{lec11:direct-sum-finite}
    \textbf{جهت عکس.}
    اگر $I$ متناهی و هر $M_i$ نوتری \textup{(آرتینی)} باشد، از
    \[
        \frac{M_1\oplus M_2}{M_1}\cong M_2
    \]
    و قضیهٔ زیرمدول و خارج‌قسمت نتیجه می‌شود $M_1\oplus M_2$
    نوتری \textup{(آرتینی)} است.

    با تکرار این استدلال،
    \[
        \bigoplus_{i=1}^{k}M_i
        =\left(\bigoplus_{i=1}^{k-1}M_i\right)\oplus M_k
    \]
    نیز نوتری \textup{(آرتینی)} است.
    \hfill\(\qed\)
\end{frame}

% Source: sources/lectures/11 Notherian and Artinian modules (2).pdf | PDF page: 6 | item: finite sum of submodules criterion
\begin{frame}[t]{مجموع متناهی زیرمدول‌ها}
    \label{lec11:finite-sum-criterion}
    \begin{theorem}
        فرض کنید $M=\sum_{i\in I}M_i$، که $I$ متناهی و هر $M_i$
        زیرمدولی از $M$ است. آنگاه $M$ نوتری \textup{(آرتینی)} است
        اگر و تنها اگر هر $M_i$ نوتری \textup{(آرتینی)} باشد.
    \end{theorem}
    \textbf{جهت مستقیم.}
    چون $M_i\leqslant M$، از نوتری \textup{(آرتینی)} بودنِ $M$،
    نوتری \textup{(آرتینی)} بودنِ هر $M_i$ نتیجه می‌شود.
\end{frame}

% Source: sources/lectures/11 Notherian and Artinian modules (2).pdf | PDF pages: 6-7 | item: epimorphism from external direct sum onto submodule sum
\begin{frame}[t]{همریختی از جمع مستقیم به مجموع}
    \label{lec11:finite-sum-epimorphism}
    \textbf{جهت عکس.}
    اگر هر $M_i$ نوتری \textup{(آرتینی)} باشد، نگاشت
    \[
        \begin{aligned}
            \varphi:\bigoplus_{i=1}^{n}M_i&\longrightarrow\sum_{i=1}^{n}M_i,\\
            (m_1,\ldots,m_n)&\longmapsto m_1+\cdots+m_n
        \end{aligned}
    \]
    یک همریختی پوشاست. بنابراین
    \[
        \frac{\bigoplus_{i=1}^{n}M_i}{\ker\varphi}
        \cong\sum_{i=1}^{n}M_i.
    \]
    جمع مستقیم متناهی، و در نتیجه خارج‌قسمت آن، نوتری \textup{(آرتینی)} است.
    \hfill\(\qed\)
\end{frame}

% Source: sources/lectures/11 Notherian and Artinian modules (2).pdf | PDF page: 7 | item: semisimple-module three-way equivalence
\begin{frame}[t]{شرایط زنجیری برای مدول نیمه‌ساده}
    \label{lec11:semisimple-criterion}
    \begin{theorem}
        اگر $M$ یک $R$-مدول نیمه‌ساده باشد، موارد زیر معادل‌اند:
        \begin{enumerate}
            \item $M$ با تولید متناهی است.
            \item $M$ آرتینی است.
            \item $M$ نوتری است.
        \end{enumerate}
    \end{theorem}
\end{frame}

% Source: sources/lectures/11 Notherian and Artinian modules (2).pdf | PDF page: 7 | item: finite supports, explicit source exercise “why?” retained
\begin{frame}[t]{تولید متناهی در مدول نیمه‌ساده}
    \label{lec11:semisimple-finite-support}
    \textbf{شرط ۱ به شرط ۲.}
    از نیمه‌ساده بودن و تولید متناهی داریم
    \[
        M=\bigoplus_{i\in I}M_i,
        \qquad M_i\ \text{ساده},
        \qquad M=\langle\alpha_1,\ldots,\alpha_k\rangle.
    \]
    جزوه در اینجا متناهی بودنِ $I$ را نتیجه می‌گیرد و می‌پرسد: \textbf{چرا؟}

    \medskip
    \textbf{راهنمای خود جزوه:}
    هر $\alpha_j$ مجموعی متناهی از عناصر ناصفرِ تعدادی از $M_i$هاست؛
    اندیس‌های ظاهرشده را در نظر بگیرید.

    \medskip
    هر $M_i$ ساده و در نتیجه نوتری و آرتینی است. با متناهی بودنِ $I$،
    $M$ نیز نوتری و آرتینی خواهد بود.
\end{frame}

% Source: sources/lectures/11 Notherian and Artinian modules (2).pdf | PDF page: 7 | item: Artinian implies Noetherian, Noetherian implies finitely generated
\begin{frame}[t]{پایان معیار مدول نیمه‌ساده}
    \label{lec11:semisimple-end}
    \textbf{شرط ۲ به شرط ۳.}
    از
    \[
        M=\bigoplus_{i\in I}M_i\quad(M_i\ \text{ساده})
    \]
    و آرتینی بودنِ $M$، با قضیهٔ جمع مستقیم نتیجه می‌شود $I$ متناهی است.
    پس $M$ نوتری است.

    \medskip
    \textbf{شرط ۳ به شرط ۱.}
    این نتیجه با قضیهٔ معیار نوتری بودن روشن است.
    \hfill\(\qed\)
\end{frame}

% Source: sources/lectures/11 Notherian and Artinian modules (2).pdf | PDF page: 8 | item: nonzero Artinian module has simple submodule; source gives only a descending-chain sketch
\begin{frame}[t]{وجود زیرمدول ساده}
    \label{lec11:artinian-simple-submodule}
    \begin{theorem}
        هر $R$-مدول آرتینی ناصفر حداقل یک زیرمدول ساده دارد.
    \end{theorem}
    \textbf{طرح اثباتِ جزوه.}
    اگر مدول ساده نباشد، با انتخاب زیرمدول‌های ناصفرِ سره،
    زنجیره‌ای نزولی ساخته می‌شود:
    \[
        M_1\supsetneq M_2\supsetneq M_3\supsetneq\cdots.
    \]
    آرتینی بودن مانع ادامهٔ نامتناهی این فرایند است؛ زیرمدول ناصفری که
    دیگر زیرمدول ناصفرِ سره ندارد، ساده است.
    \hfill\(\qed\)
\end{frame}

% Source: sources/lectures/11 Notherian and Artinian modules (2).pdf | PDF page: 8 | item: regular-module vs all finitely generated modules, statement
% AMBIGUITY[L11-Q04]: unital left modules are an explicit provisional working assumption, not a resolved source convention.
\begin{frame}[t]{مدول‌های با تولید متناهی روی حلقه}
    \label{lec11:regular-module-criterion-pending}
    \begin{tblock}{فرض موقت: مدول‌ها چپ و یکانی‌اند؛ \lr{L11-Q04}}
        اگر $R$ حلقه‌ای یکدار باشد، موارد زیر معادل‌اند:
        \begin{enumerate}
            \item $R$ به‌عنوان $R$-مدول چپ نوتری \textup{(آرتینی)} است.
            \item هر $R$-مدول چپ یکانی با تولید متناهی نوتری \textup{(آرتینی)} است.
        \end{enumerate}
    \end{tblock}
    قید یکانی بودن، انتخاب موقتِ این پیش‌نویس است؛ در صورت دست‌نویس صریح نیست.
    این پرسش در پروندهٔ جداگانه باز باقی مانده است.
\end{frame}

% Source: sources/lectures/11 Notherian and Artinian modules (2).pdf | PDF page: 8 | item: free-module epimorphism proof and reverse via R=(1)
% AMBIGUITY[L11-Q04]: source proof under the visibly provisional unital-left-module convention.
\begin{frame}[t]{طرح اثبات با خارج‌قسمت مدول آزاد}
    \label{lec11:regular-module-proof-pending}
    \textbf{با فرض موقتِ یکانی بودن: \lr{L11-Q04}.}
    برای $M=\langle m_1,\ldots,m_n\rangle$ نگاشت
    \[
        \begin{aligned}
            \varphi:R^n&\longrightarrow M,\\
            (r_1,\ldots,r_n)&\longmapsto r_1m_1+\cdots+r_nm_n
        \end{aligned}
    \]
    یک همریختی پوشاست و
    \[
        M\cong R^n/\ker\varphi.
    \]
    از نوتری \textup{(آرتینی)} بودنِ $R$، همین خاصیت برای $R^n$
    و سپس خارج‌قسمت آن نتیجه می‌شود.

    \medskip
    برای عکس، جزوه از $R=(1)$ استفاده می‌کند.
\end{frame}

% Source: sources/lectures/11 Notherian and Artinian modules (2).pdf | PDF page: 8 | item: simple submodule over left Artinian ring, final theorem and proof
% AMBIGUITY[L11-Q04]: explicit unital-left-module convention is a provisional working assumption; source says arbitrary nonzero R-module.
\begin{frame}[t]{زیرمدول ساده بر حلقهٔ آرتینی}
    \label{lec11:left-artinian-simple-pending}
    \begin{tblock}{با فرض موقتِ یکانی بودن: \lr{L11-Q04}}
        اگر $R$ یکدار و به‌عنوان $R$-مدول چپ آرتینی باشد، هر $R$-مدول چپ
        یکانیِ ناصفر $M$ حداقل یک زیرمدول ساده دارد.
    \end{tblock}
    \textbf{طرح اثباتِ جزوه با این فرض موقت.}
    برای $0\ne x\in M$، زیرمدول $N=Rx$ ناصفر و با تولید متناهی است.
    پس $N$ آرتینی است و بنابر قضیهٔ قبل یک زیرمدول ساده دارد.
    \hfill\(\qed\)

    \medskip
    یکانی بودن در دست‌نویس تصریح نشده؛ انتخاب موقت در پروندهٔ پرسش‌ها ثبت شده است.
\end{frame}
